\documentclass[UTF8,a4paper,11pt]{ctexart}
\usepackage{amsmath,amssymb}
\usepackage[margin=2.5cm]{geometry}
\usepackage[hidelinks]{hyperref}
\title{2025年上海卷第19题：含参对数函数的极大值}
\date{2026年10月24日}
\begin{document}
\maketitle
\noindent\textbf{固定来源：}\url{https://github.com/DxAThing/Gaokao-Math-Problems-Compilation/blob/f1ae169b58a75b3db6e35c670b3f0d5c78718188/content/2025/shanghai.tex#L381}

设
\[
 f_m(x)=x^2-(m+2)x+m\ln x,\quad x>0,\ m\in\mathbb R.
\]
\paragraph{一、求不等式解集}
$f_m(1)=-m-1=0$ 导致 $m=-1$。
此时 $f_{-1}(x)\le x^2-1$ 等价于
$x+\ln x\ge1$。因为
\[
\frac{\mathrm d}{\mathrm dx}(x+\ln x)=1+\frac1x>0
\quad(x>0)
\]
且 $1+\ln1=1$，故解集为
\[
\boxed{[1,+\infty)}.
\]

\paragraph{二、极大值存在的参数}
求导并分解为
\[
f_m'(x)=\frac{(2x-m)(x-1)}{x}.
\]
\begin{enumerate}
\item $m\le0$ 时，合法驻点仅 $x=1$，
导数由负变正，仅有极小值。
\item $0<m<2$ 时，$0<m/2<1$，
导数正、负、正，在 $x=m/2$ 取得严格极大值。
\item $m=2$ 时，$f_m'(x)=2(x-1)^2/x\ge0$，
函数严格递增，没有极大值。
\item $m>2$ 时，$1<m/2$，导数正、负、正，
在 $x=1$ 取得严格极大值。
\end{enumerate}
故参数范围为
\[
\boxed{m\in(0,2)\cup(2,+\infty)}.
\]
特别注意 $m=2$ 的重驻点不是极值点。

\bigskip
\noindent\textbf{审校状态：}数学独立复核；
新源码尚未实际编译。
\end{document}
